% answers_short.tex —— 2006 年·普通高等学校招生考试·陕西卷文科 · 简短答案速查
% 由 tools/make_short_answers.py 自 answers.tex 的 \gksolhead/\gkanswer 提取，
% 并与 2006/2006_数学真题_陕西文_答案要点.md 逐题比对后生成，勿手改（改源 MD 后重跑）。

\begin{gknotice}
\textbf{使用说明}：本节只给最终答案便于快速核对，完整推导见解析版。答案为编写者独立推导，\textbf{非官方标准答案}。
\end{gknotice}

\begin{gknotice}
选择题答案中的 U 表示待核，未给出猜测答案，详情见解析版。
\end{gknotice}

\gkpart{选择题}
\begin{center}\begin{tabular}{ccccccccccccc}
\toprule
题号 & 1 & 2 & 3 & 4 & 5 & 6 & 7 & 8 & 9 & 10 & 11 & 12 \\
\midrule
答案 & A & B & U & C & B & A & B & D & A & D & D & D \\
\bottomrule
\end{tabular}\end{center}

\gkpart{填空题}
\begin{itemize}[leftmargin=2.6em]
  \item[\textbf{13.}] $-$1/2。
  \item[\textbf{14.}] 60。
  \item[\textbf{15.}] U/PENDING。
  \item[\textbf{16.}] 3R。
\end{itemize}

\gkpart{解答题}
\begin{itemize}[leftmargin=2.6em]
  \item[\textbf{17.}] (1) 3/25；(2) 19/50
  \item[\textbf{18.}] (1) 最小正周期 $\pi$；(2) $x=\frac{5\pi}{12}+k\pi$，$k\in\mathbb{Z}$
  \item[\textbf{19.}] (1) 直线 $AB$ 与 $\alpha$、$\beta$ 的夹角依次为 $\frac\pi4$、$\frac\pi6$；(2) $\arccos(1/\sqrt3)$
  \item[\textbf{20.}] $a_n=5n-3$
  \item[\textbf{21.}] (1) 动直线斜率范围为 $[-1,1]$；(2) $y=\frac{x^2}{4}$，$-2\le x\le2$
  \item[\textbf{22.}] (1) $k=0$ 时，$f$ 在 $(-\infty,0]$ 上递增、在 $[0,+\infty)$ 上递减；$k>0$ 时，$f$ 在 $(-\infty,0]$ 与 $[2/k,+\infty)$ 上递增，在 $[0,2/k]$ 上递减；(2) $k>2$
\end{itemize}
